%%%%%%%%%%%%%%%%%%%%%%%%%%%%%%%%%%%%%%%%%
% Journal Article
% LaTeX Template
% Version 2.0 (February 7, 2023)
%
% This template originates from:
% https://www.LaTeXTemplates.com
%
% Author:
% Vel (vel@latextemplates.com)
%
% License:
% CC BY-NC-SA 4.0 (https://creativecommons.org/licenses/by-nc-sa/4.0/)
%
% NOTE: The bibliography needs to be compiled using the biber engine.
%
%%%%%%%%%%%%%%%%%%%%%%%%%%%%%%%%%%%%%%%%%

%----------------------------------------------------------------------------------------
%	PACKAGES AND OTHER DOCUMENT CONFIGURATIONS
%----------------------------------------------------------------------------------------

\documentclass[
	a4paper, % Paper size, use either a4paper or letterpaper
	10pt, % Default font size, can also use 11pt or 12pt, although this is not recommended
	unnumberedsections, % Comment to enable section numbering
	twoside, % Two side traditional mode where headers and footers change between odd and even pages, comment this option to make them fixed
]{LTJournalArticle}

\addbibresource{bibliography.bib} % BibLaTeX bibliography file

\runninghead{Shortened Running Article Title} % A shortened article title to appear in the running head, leave this command empty for no running head

\footertext{\textit{Journal of Biological Sampling} (2024) 12:533-684} % Text to appear in the footer, leave this command empty for no footer text

\setcounter{page}{1} % The page number of the first page, set this to a higher number if the article is to be part of an issue or larger work
\emergencystretch=2em % Reduce underfull/overfull hbox warnings by relaxing line breaking

%----------------------------------------------------------------------------------------
%	TITLE SECTION
%----------------------------------------------------------------------------------------

\title{Simulation et modélisation de  trajectoires circulaires} % Article title, use manual lines breaks (\\) to beautify the layout

% Authors are listed in a comma-separated list with superscript numbers indicating affiliations
% \thanks{} is used for any text that should be placed in a footnote on the first page, such as the corresponding author's email, journal acceptance dates, a copyright/license notice, keywords, etc
\author{%
	Manuel Rocca, Damian Kazberuk, Matteo Morbée
}

% Affiliations are output in the \date{} command
\date{\today}

% Full-width abstract
\renewcommand{\maketitlehookd}{%
	\begin{abstract}
		\noindent ABSTRACT
	\end{abstract}
}

%----------------------------------------------------------------------------------------

\begin{document}

\maketitle % Output the title section

%----------------------------------------------------------------------------------------
%	ARTICLE CONTENTS
%----------------------------------------------------------------------------------------

\section{Introduction}
Quel est un bon modèle pour représenter un phénomène réel ? Plus particulièrement, comment choisir un modèle pour un cas précis ? Pour répondre à ces questions, nous avons élaboré plusieurs modèles sur base de notre expérience physique, la trajectoire d'une bille tournant sur une surface "placer lien vers figure expérience ?". Ce rapport présente et analyse quatres modèles, deux modèles basés sur des équations physiques et deux modèles basés sur des méthodes de machine learning.

%------------------------------------------------

\section{Modèles}

\subsection{Modèle conique}
\subsubsection{Équations}
\subsubsection{Limitations et hypothèses}

\subsection{Modèle logarithmique}
\subsubsection{Équations}
\subsubsection{Limitations et hypothèses}

\subsection{Modèle de régression linéaire}
\subsubsection{Équations}
\subsubsection{Limitations et hypothèses}

\subsection{Modèle de "multilayer perceptron" MLP}
\subsubsection{Équations}
\subsubsection{Limitations et hypothèses}

%------------------------------------------------

\section{Simulation}


\subsection{Simulation physique}
Première grande partie de notre réflexion, la simulation physique nous a permis de modéliser la trajectoire de la bille à l'aide de différentes méthodes numériques. Nous avons utilisé des méthodes d'Euler explicite, semi-implicite et Velvet, ainsi que la méthode de Runge-Kutta d'ordre 4 (RK4) pour simuler le mouvement de la bille sur la surface circulaire. Ces méthodes nous ont permis d'obtenir des résultats précis et de comparer les différentes approches pour simuler le même phénomène. Nous avons également manipulé les paramètres de notre modèle pour observer comment ils affectent la trajectoire de la bille, ce qui nous a permis de mieux comprendre les forces en jeu et les facteurs qui influencent le mouvement de la bille. Ces observations nous ont permis de tirer des conclusions sur la précision et la fiabilité de nos simulations, ainsi que sur les limites de notre modèle physique. En effet, nous avons constaté que représenter la réalité de manière précise relève de l'exploit, et que, en plus des approximations numériques inévitables, il est nécessaire de faire des hypothèses et des simplifications pour rendre le modèle plus facile à comprendre et à résoudre. Par exemple, nous avons fait  plusieurs hypothèse: la pression de l'air n'impact pas la trajectoire de la bille, la déformation de la surface par la bille est négligeable, la bille est un point matériel, etc. Nous avons donc réalisé plusieurs expériences sur deux représentations différentes de notre surface circulaire, une utilisant une forme géométrique conique et l'autre utilisant une forme logarithmique. Cela couplé avec les différentes méthodes d'intégration numérique et à la manipulation (ajout et retrait) des paramètres nous a permis d'obtenir une grande variété de résultats et d'observations. % ajouter graphique ici ??


\subsection{Simulation à l'aide de machine learning}
La deuxième grande famille de méthodes de simulation, est celle utilisant le machine learning. 
Celles-ci fonctionnant avec une grande quantité de données, nous avons utilisé les résultats obtenus à partir de nos 
simulations physiques pour entraîner nos modèles de machine learning.
Nous avons utilisé deux types de modèles de machine learning: la régression linéaire et le multi-layer perceptron (MLP).
La régression linéaire est un modèle de machine learning simple qui permet de modéliser la relation entre une variable dépendante 
et une ou plusieurs variables indépendantes. 
Nous avons utilisé la régression linéaire pour modéliser la relation entre les paramètres de notre modèle physique, tels que le vecteur 
vitesse initial, ainsi que la position de la bile a chaque instant t, et les résultats obtenus est donc la trajectoire de la bille. 
Cette approche nous a permis d'obtenir une approximation de plus en plus précise en fonction du nombre de données disponibles.
Ceci nous a permis de reproduire la trajectoire de la bille à partir de données sans avoir la connaissance ni la compréhension du modèle physique sous-jacent.

Le multi-layer perceptron (MLP) est un modèle de machine learning plus complexe qui permet de modéliser des relations non linéaires entre les variables. 
Nous avons utilisé le MLP pour modéliser la relation entre les paramètres de notre modèle physique et les résultats obtenus, de manière plus précise que la régression linéaire. 
Cette approche nous a permis d'obtenir des résultats plus précis que la régression linéaire, en particulier pour les trajectoires plus complexes. 
Cependant, le MLP nécessite une plus grande quantité de données pour être entraîné de manière efficace.

Cette approche de simulation à l'aide de machine learning nous a permis d'obtenir des résultats probant et réaliste, sa force résidant 
dans sa capacité à apprendre à partir de données et à modéliser des relations complexes entre les variables fait du machine learning
 une approche très intéressante pour simuler des phénomènes complexes, 
 en particulier lorsque la compréhension du modèle physique sous-jacent est limitée. Cependant, il est important de noter que le machine learning 
 ne remplace pas la compréhension du modèle physique, et que les résultats obtenus à partir de modèles de machine learning doivent être 
 interprétés avec prudence, en tenant compte des limites et des hypothèses du modèle utilisé.
%------------------------------------------------

\section{Calcul numérique}

\subsection{Approximation numérique}

Dans le cadre d'une approche scientifique, il est important de prendre en compte tous les paramètres pouvant entrer en jeu et potentiellement perturber et modifier l'expérience. Dans notre cas nous distinguons deux types d'erreur principaux\cite[16]{Marteen:2022}: les erreurs numériques et les erreurs de modélisation.

\subsubsection{Erreurs numériques}

Les erreurs numériques apparaissent lorsque l'ordinateur effectue des calculs et se déclinent sous différentes formes.
Dans un premier temps, les erreurs de troncature sont des erreurs qui apparaissent lorsque nous devons arrêter un processus qui est infini, par exemple une série de Taylor.

Ensuite, les erreurs d'arrondi sont des erreurs causées par la représentation finie des nombres en machine dans les registres du processeur (généralement 32 ou 64 bits). L'ordinateur ne peut pas représenter tous les nombres réels avec une précision infinie, ce qui crée alors une approximation.

Les erreurs de propagation et de génération sont des erreurs qui apparaissent lorsque les erreurs numériques se propagent à travers les calculs. Les erreurs de génération se produisent à chaque itération, étape des calculs et les erreur de propagation regroupent toutes les erreurs de génération jusqu'à l'étape courante. Par exemple, si nous simplifions $\pi$ à 3.14, nous introduisons une erreur d'arrondi à une certaine génération qui va se propager à travers les calculs ultérieurs, affectant ainsi la précision globale des résultats.

\subsubsection{Erreurs de modélisation}

La catégorie des erreurs de modélisation se décline en deux types d'erreurs.

Le premier type, les erreurs de mesure qui se produisent dans le cadre physique de notre expérience. En effet, nous ne disposons pas d'outils de mesure parfaits, ce qui induit donc des imprécisions numériques de l'ordre du millimètre lorsque nous mesurons des éléments de notre expérience comme par exemple la largeur du dispositif.

Le second regroupe les erreurs de modélisation du phénomène; des erreurs qui apparaissent lorsque nous faisons des simplifications, des hypothèses dans notre modèle pour le rendre plus facile à comprendre ou à résoudre. Un exemple d'hypothèse émise dans le cadre de notre expérience concerne la pression de l'air, supposée sans impact sur la trajectoire de la balle.

\subsection{Consistance}
La consistance est la capacité d'un algorithme à produire des résultats qui convergent vers la solution exacte du problème.
Un algorithme est dit consistant si, à mesure que la discrétisation devient plus fine, les résultats obtenus se rapprochent de la solution exacte.
Dans notre cas, la méthode de Euler (explicite ou implicite) est consistante. % d'ordre 1. Parler des autres aussi quand on aura fait/trouver les étapes de calculs.

\subsection{Stabilité}
La stabilité est la propriété d'un algorithme qui garantit que les erreurs apportées durant les calculs ne s'amplifient pas de manière exponentielle au fil des itérations.
Un algorithme est dit stable si une petite perturbation dans les données d'entrée entraîne une petit perturbation dans les résultats. % meh mais dans le cours mdr

\subsection{Conditionnement}
Le conditionnement décrit la sensibilité intrinsèque du problème : une petite perturbation des données d'entrée peut produire une petite ou une grande variation de la solution exacte.
Un problème bien conditionné est peu sensible aux perturbations, tandis qu'un problème mal conditionné les amplifie.
% TODO: ASPECT ANALYTIQUE
% en particulier, il nous faut un cas concret, une erreur concrète, prise en flagrant délit pour appuyer nos propos
% => comment faire ?: outils de mesure par exemple MSE (ou tout autre)

% Parler du conditionnement du problème ? consistance ?

% notation des nombres en machine, précision finie => EN PARLER ABSOLUMENT (tout simplement car on en a beaucoup parlé dans la vulgarisation et car c'est un peu la source de tous les problèmes en qlq sorte)


% différentes simus:
% simus physique: euler explicite, semi-implicite et velvet (+rk4)
% simu ml: régression linéaire et MLP (multi-layer perceptron)
%------------------------------------------------

\section{Génération et analyse des données}

\subsection{Pipeline de génération}
Pour produire les données nécessaires à l'analyse, nous avons exécuté un pipeline reproductible en quatre étapes:
\begin{enumerate}[leftmargin=*, itemsep=0.2em, topsep=0.2em]
	\item génération des trajectoires synthétiques
	\item entraînement des modèles
	\item exécution des benchmarks
	\item comparaison avec les données réelles
\end{enumerate}

La génération synthétique a simulé \(1\,000\,000\) trajectoires, dont \(375\,266\) conservées après filtrage, pour un total de \(103\,792\,281\) paires d'états utilisées pour l'entraînement.

\subsection{Résultats quantitatifs clés}
\begin{table}[h]
	\centering
	\small
	\setlength{\tabcolsep}{4pt}
	\caption{Indicateurs extraits des résultats de benchmark.}
	\begin{tabularx}{\linewidth}{|>{\raggedright\arraybackslash}p{0.36\linewidth}|>{\raggedright\arraybackslash}X|}
		\hline
		\textbf{Indicateur} & \textbf{Valeur observée} \\
		\hline
		Convergence intégrateurs (pente log-log) & Euler explicite \(\approx 1.20\), Euler semi-implicite \(\approx 1.11\), RK4 \(\approx 4.00\) \\
		\hline
		RMSE radial à \(\Delta t=0.01\,s\) & Euler explicite \(9.23\times 10^{-4}\,m\), Euler semi-implicite \(6.47\times 10^{-4}\,m\), RK4 \(1.93\times 10^{-9}\,m\) \\
		\hline
		Meilleur modèle ML step-by-step & MLP (10\%) : \(\mathrm{MAE}_r = 3.55\times 10^{-3}\,m\), stabilité \(71.5\%\) \\
		\hline
		Comparaison sur test réel & MAE radiale normalisée : Linéaire \(0.109\,R\), MLP \(0.217\,R\) \\
		\hline
		Erreur récursive à 5 s (\(|\Delta r|/R\), médiane) & Linéaire \(0.0597\), MLP \(0.0525\) \\
		\hline
	\end{tabularx}
\end{table}

\subsection{Lecture des résultats}
Les mesures physiques confirment la théorie: RK4 est nettement plus précis, mais plus coûteux; Euler semi-implicite reste un compromis robuste pour des pas de temps usuels. Côté ML, les performances sur données synthétiques sont bonnes, mais la comparaison au réel montre un écart significatif (\textit{sim-to-real}).

L'analyse d'horizon confirme également un point important: même quand l'erreur au pas est faible, la prédiction récursive accumule les écarts. Ce résultat justifie l'usage de métriques dépendant du temps (et pas seulement d'une erreur moyenne globale) pour évaluer les modèles dynamiques.

\section{Discussion et conclusion}

%----------------------------------------------------------------------------------------
%	 ANNEXES
%----------------------------------------------------------------------------------------
% On peut mettre le développement de la consistance des méthodes ? Euler dans le cours de CFN

\printbibliography % Output the bibliography

%----------------------------------------------------------------------------------------

\end{document}
